\begin{helper}
Let $G$ be a finite simple graph and let $v$ be a vertex.  Form the graph
$C$ on the neighbor set $N_G(v)$ by joining two distinct neighbors exactly
when they are not adjacent in $G$.  Then
\[
|E(C)|\le P_G(v),
\]
where $P_G(v)$ is the number of independent pairs in $N_G(v)$ as defined in
\noderef{IndependentPairCount}.
\end{helper}
\begin{proof}
Each edge of $C$ is an unordered pair $\{x,y\}$ of distinct vertices in
$N_G(v)$ such that $x$ and $y$ are nonadjacent in $G$.  Sending the edge to
the two-element subset $\{x,y\}$ therefore gives an independent pair in the
neighborhood of $v$, exactly of the kind counted by
\noderef{IndependentPairCount}.  The map is injective, because a simple-graph
edge is determined by its two endpoints as an unordered pair.  Hence the
number of edges of $C$ is at most the number of independent neighbor-pairs
counted by $P_G(v)$.
\end{proof}
